Grandes Modelos de Linguagem (LLMs) são utilizados na geração de conteúdo e no apoio a decisões. O alinhamento por preferências pode criar tensões entre produzir respostas apreciadas pelo usuário e preservar a correção factual. Nesse contexto, \textit{sycophancy} designa comportamentos de concordância indevida com crenças ou sugestões do usuário. Este trabalho propõe avaliar, como desfecho principal, a capitulação factual: adoção da alternativa falsa sugerida pelo usuário após resposta inicialmente correta. O protocolo candidato utiliza perguntas de múltipla escolha com gabaritos verificáveis e respostas iniciais reais. Após acerto, serão comparados ramos de sugestão falsa sem nova evidência e seguimento neutro; após erro real, a aceitação de correção verdadeira será examinada como análise complementar, também com controle neutro. Propõem-se dois seguimentos por ramo, repetindo a mesma mensagem e preservando o histórico, sem estimar um efeito isolado de intensidade. As medidas incluem capitulação específica, degradação factual geral, aceitação de correção verdadeira, primeiro turno do evento, estado final e recuperação. Serão explicitadas as populações elegíveis, a cobertura e as respostas indeterminadas, com classificação automática, validação humana e análise que considere a dependência entre repetições da mesma questão. A criticidade será exploratória, distinta de domínio e dificuldade, representando consequências potenciais dos cenários, sem medir dano real ou segurança geral dos modelos. O protocolo permanece proposto, sujeito à avaliação do grupo e do orientador e à revisão após o piloto; não são apresentados resultados experimentais de LLMs.
